\chapter{Introducción}

En esta sección se presenta la motivación para la elaboración de este proyecto, los objetivos a alcanzar, la metodología seguida y la estructura del documento.

\section{Motivación}

En el ámbito de la ciberseguridad, los sistemas de detección de intrusiones (\textit{Intrusion Detection Systems}, IDS) desempeñan un papel fundamental 
en la identificación temprana de ataques y actividades maliciosas en redes informáticas. En los últimos años, el uso de técnicas de inteligencia artificial 
y aprendizaje automático (\textit{Machine Learning}) ha permitido mejorar significativamente la precisión y la capacidad de adaptación de estos sistemas 
frente a amenazas cada vez más complejas y dinámicas.

Sin embargo, gran parte de la investigación existente se ha centrado principalmente en mejorar la tasa de detección o reducir 
los falsos positivos, sin considerar en la misma medida el impacto computacional que conllevan estos modelos. En entornos 
donde el procesamiento debe realizarse en tiempo real, como redes corporativas, sistemas embebidos, dispositivos IoT o 
infraestructuras críticas, el consumo de recursos, la memoria utilizada y la latencia de inferencia se convierten en factores 
determinantes para la viabilidad de la solución.

Por este motivo, resulta necesario analizar no solo la capacidad predictiva de los modelos, sino también su eficiencia y su comportamiento ante escenarios adversos. 
En este contexto, este Trabajo Fin de Grado plantea una evaluación experimental de distintos modelos de aprendizaje automático y aprendizaje profundo aplicados a la detección de intrusiones, considerando tanto su rendimiento en condiciones normales como su robustez frente a perturbaciones adversarias.

\section{Objetivos}

El objetivo principal de este Trabajo de Fin de Grado es evaluar de forma integral el rendimiento, la robustez y la viabilidad operativa de diferentes modelos de aprendizaje automático y profundo aplicados a Sistemas de Detección de Intrusiones (IDS). Se busca analizar el compromiso (trade-off) existente entre la capacidad predictiva de los algoritmos frente a amenazas camufladas y el coste computacional requerido para su ejecución en entornos con recursos limitados.

\section{Metodología de trabajo}

Para la elaboración de este TFG, se ha seguido una metodología estructurada siguiendo un enfoque iterativo y experimental. Dada la naturaleza exploratoria del proyecto, se ha seguido una metodología mista, combinando el modelo CRISP-DM con un enfoque experimental orientado al análisis de rendimiento.
Esta metodología se divide en las siguientes fases:

\begin{itemize}
    \item Estudio del estado del arte sobre aprendizaje automático y sistemas de detección de intrusiones.
    \item Obtención y análisis de los datasets, evaluando su viabilidad para el proyecto.
    \item Implementación, evaluación y búsqueda de hiperparámetros de modelos de Machine Learning y Deep Learning buscando un equilibrio óptimo entre detección y uso de recursos.
    \item Aplicación de ataques de evasión a los modelos ganadores
    \item Análisis y conclusiones de los resultados obtenidos 
\end{itemize}

Además, se ha tenido un contacto continuo con los tutores, a través de mensajería y reuniones frecuentes para seguir el proceso continuo de avance del TFG y plantear las dudas correspondientes.

\section{Estructura del documento}

Esta memoria está estructurada en los siguientes capítulos:

\begin{itemize}
    \item \textbf{Capítulo 1. Introducción: } se presenta una introducción del proyecto, a través de su motivación, los objetivos y la metodología de trabajo a seguir.
    \item \textbf{Capítulo 2. Investigación previa: } se tratan los conceptos básicos para entender el objetivo del proyecto. Además, se explica el contexto y estado actual de la detección de anomalías a través de aprendizaje automático.
    \item \textbf{Capítulo 3. Tecnologías utilizadas: } se describen las tecnologías utilizadas por el proyecto. 
    \item \textbf{Capítulo 4. Obtención y Análisis de Datos: } se detallan los datasets elegidos con un análsis de ellos.
    \item \textbf{Capítulo 5. Análisis de Modelos y Búsqueda de Hiperparámetros: } se explica el proceso seguido para la búsqueda de los mejores hiperparámetros de cada modelo, junto con sus métricas y resultados.
    \item \textbf{Capítulo 6. Ataques de Evasión: } se tratan los modelos elegidos y se exponen ante ataques adversarios para ver la degradación de la eficacia de los modelos.
    \item \textbf{Capítulo 7. Conclusiones y líneas futuras: } se reúnen conclusiones obtenidas tras la realización del trabajo, además de caminos de interés a seguir para trabajos futuros.
\end{itemize}